% page 172 du fac-simile (image p-171 du scan)
% bloc 165.1 x 246.3 mm ; marge gauche 8.0 mm
\pgc{-12.700mm}{0.000mm}{
\l{\cel{3}162}
\l{}
\l{}
\l{\sou{fenolo}. (\sou{kemio}).Substanco desinfektiva, extraktita de la olei}
\l{\cel{3}tre pezoza, devenanta de la distilo di minkarbono e di gu-}
\l{\cel{3}dro.\cel{1}\sou{Simbolo kemiala}\cel{1}: C6H5\cel{2}OH. - DEFIRS.}
\l{}
\l{\sou{fenomeno}. Modifikeso di korpo quan perceptas nia sensi.-DEFIRS.}
\l{}
\l{\sou{fenomenologio.}\cel{1}I. Traktato, diserturo, pri la fenomeni, pri}
\l{\cel{3}to quo povas frapar la sensi. - II. (\sou{filoz}.)En la Hegel-}
\l{\cel{3}sistemo, cienco pri la idei qui naskas de la percepti da la}
\l{\cel{3}sensi. - DEF.}
\l{}
\l{\sou{fero}. Metalo blanka, grizatra, fuzebla per alta temperaturo,}
\l{\cel{3}tenaca,duktila, forjebla, magnetala,uzata tre diverse en in-}
\l{\cel{3}dustrio. -\cel{1}\sou{Simbolo kemiala}\cel{1}: Fe. - EFIS.}
\l{}
\l{\sou{ferdeko}. (\sou{nav}o).Plankaro fixigita sinse di la longeso di navo,}
\l{\cel{3}por kovrar la holdo o por formacar etaji. - DE.}
\l{}
\l{\sou{ferio}. Granda merkato ube vendesas vari omnasorta, qua eventas}
\l{\cel{3}unfoye o plurafoye singlayare, en regiono, ordinare lor ed}
\l{\cel{3}en la loko ube celebresas la festo lokala. - EF.}
\l{}
\l{\sou{ferma}. Qua ne flexas, qua ne shancelas.\cel{3}II. (\sou{metaf}.) Qua ne}
\l{\cel{3}ocilas. - III. (\sou{komerco)}. Di qua la kondicioni esas defini-}
\l{\cel{3}tiva e ne dependas de ta o ca eventualajo. - DEFIS.}
\l{}
\l{\sou{fermento}. I. To quo produktas ula transformi molekulala (dia-}
\l{\cel{3}stazi, e c.)- II. (\sou{metaf.}) To quo naskigas ula agitesi psi-}
\l{\cel{3}kala. - DEFIRS.}
\l{}
\l{\sou{fermentacar}. (\sou{netrans.}) Subisar la transformo molekulala quan}
\l{\cel{3}produktas fermento. (\sou{anke metaf.}) - DEFIRS.}
\l{}
\l{\sou{feroca}.\cel{2}Di qua la karaktero esas kruela, sango-durstanta, -EFIS.}
\l{}
\l{\sou{fertila}. Qua produktas multo (exemple : agro). - DEFIS.}
\l{}
\l{\sou{fervorar}. (\sou{netrans., pri, pro, por}). Qua ardoras.vivace. - EFIS.}
\l{}
\l{\sou{\textquotedbl{}fes\textquotedbl{}}. Noto quaresma di la \textquotedbl{}do\textquotedbl{}gamo, bemoloza.}
\l{}
\l{\sou{festar}. I. (\sou{trans.})\cel{2}Agar la solenaji religiala, memorigala,}
\l{\cel{3}quin onu celebras en ula dii di la yaro ; joyo-manifestar}
\l{\cel{3}honorize di eventajo memorinda. - II. (\sou{netrans.}) Partoprenar}
\l{\cel{3}gala-festino, gala-balo, e c., donata da rejo, princo, urbo,}
\l{\cel{3}privato. - DEFIS.}
\l{}
\l{\sou{festinar.}\cel{2}(\sou{netrans.)}\cel{2}Agar repasto festala, pompoza, honorize}
\l{\cel{3}di ulo o di ulu. - EFS.}
\l{}
\l{festono. I. Envolvajo ek folii e flori, arkoforma, quan onu}
\l{\cel{3}suependas intervalope, kom festo-ornivi. - II. Bordo di vesto,}
\l{\cel{4}di linjo, taliita formacante denti kurva (o quadrata, o tri-}
\l{\cel{3}angula). -III. Broduro qua kovras la bordo. - DEFIRS.}
}
